\documentclass[11pt]{article}

\usepackage[a4paper, margin=2.6cm]{geometry}
\usepackage{amsmath, amssymb, amsthm}
\usepackage{bm}
\usepackage{graphicx}
\usepackage{booktabs}
\usepackage{siunitx}
\usepackage[dvipsnames]{xcolor}
\usepackage[colorlinks=true, linkcolor=MidnightBlue, citecolor=MidnightBlue, urlcolor=MidnightBlue]{hyperref}
\usepackage[round]{natbib}
\usepackage{microtype}

% ---- symbols ----
\newcommand{\R}{\mathbb{R}}
\newcommand{\dd}{\mathrm{d}}
\newcommand{\state}{u}                 % transported field
\newcommand{\lat}{\zeta}               % latent regime coordinate
\newcommand{\barrier}{\beta}           % barrier coordinate in [0,1]
\newcommand{\drive}{a}                 % normal-form drive
\newcommand{\ctrl}{\bm{w}}             % exogenous inputs
\newcommand{\emiss}{y}                 % scalar emission observable
\newcommand{\pH}{p_{\mathrm{H}}}

\title{Identification of a bistable normal form for regime transitions\\
in a driven dissipative transport system}
\author{Kyounghyun Lee}
\date{\today}

\begin{document}
\maketitle

\begin{abstract}
Spatially extended dissipative systems under strong forcing can switch abruptly
between distinct macroscopic transport regimes. We propose a hybrid
reduced-order model in which a one-dimensional diffusion equation with a
state-dependent transport coefficient is coupled to a scalar latent variable
governed by the cusp normal form
$\tau\,\dot\lat = \drive(\ctrl) + b\,\lat - \lat^{3}$.
The two stable branches of the cusp define the two transport regimes; the
transition and its hysteresis are saddle-node bifurcations of the latent
dynamics, and the distance of the exogenous drive from the fold provides a
quantitative, closed-form margin to transition. All model components --- the
transport coefficient, a residual source term, the drive map, and an
observation map from the latent state to a measured emission proxy --- are
identified end-to-end from experimental time series by differentiating through
an implicit--explicit integrator. Unlike proxy-following approaches, the
latent dynamics are driven exclusively by exogenous actuation, so the inferred
regime state is falsifiable against held-out observables. The identified model
assigns a calibrated regime probability at every instant, locates transition
times without supervision beyond weak segment labels, and yields
bifurcation diagrams that summarize when a transition can be forced or undone.
\end{abstract}

\section{Introduction}
\label{sec:intro}

% Placeholder: motivation from dynamical systems viewpoint.
Abrupt transitions between metastable states are a recurring motif in
nonequilibrium systems: bistability, hysteresis, and critical slowing down
near folds are generic features predicted by normal-form theory
\citep{kuznetsov2004elements, strogatz2018nonlinear}. In the experimental
system considered here, a driven dissipative transport process switches
between a low-transport-barrier regime (\emph{state L}) and a
high-transport-barrier regime (\emph{state H}) as the external drive
increases, with the reverse switch occurring at a lower drive --- the
signature of hysteresis. A quantitative, model-based answer to the question
\emph{``which regime is the system in, and how far is it from switching?''}
is a prerequisite for control.

Our contributions are:
(i) a hybrid mechanistic--statistical model in which the regime is a latent
coordinate of a cusp normal form, giving the L/H classification a precise
mathematical meaning (basin membership of a bistable vector field);
(ii) end-to-end identification through a differentiable implicit--explicit
solver; and
(iii) closed-form bifurcation diagnostics --- fold margins
$\mu_{\mathrm{LH}}(t)$ and $\mu_{\mathrm{HL}}(t)$ --- that quantify
proximity to forced and spontaneous transitions.

\section{Model}
\label{sec:model}

\subsection{Transport backbone}
Let $\state(\rho, t)$ denote the transported scalar field on the normalized
domain $\rho \in [0, 1]$ with Neumann condition at $\rho = 0$ and a measured
Dirichlet trace at $\rho = 1$. The field obeys a conservative diffusion law
with a learned residual source $s_{\theta}$,
\begin{equation}
\partial_t \state
= \frac{1}{V'(\rho)} \, \partial_\rho\!\left( V'(\rho)\, \chi(\rho, \lat)\,
\partial_\rho \state \right) + s_{\theta}\!\left(\rho, \state, n, \ctrl, \lat\right),
\label{eq:pde}
\end{equation}
where $V'$ is a fixed geometric weight and $n$ an auxiliary measured field.
The transport coefficient is a monotone boundary-layer profile whose edge
value is modulated by the latent regime coordinate $\lat$,
\begin{equation}
\chi(\rho, \lat) = \chi_0 + \sigma\!\left(\tfrac{\rho - \rho_c}{w}\right)
\left[\chi_1 - \Delta\chi\,\barrier(\lat) - \chi_0\right],
\qquad
\barrier(\lat) = \sigma\!\left(\tfrac{3\lat}{\sqrt{b}}\right),
\label{eq:chi}
\end{equation}
so that the H regime ($\lat > 0$) suppresses transport in the boundary layer
and steepens the stationary profile there.

\subsection{Latent regime dynamics: the cusp normal form}
The regime coordinate evolves according to
\begin{equation}
\tau\,\dot\lat = \drive(\ctrl) + b\,\lat - \lat^{3},
\qquad b > 0,
\label{eq:cusp}
\end{equation}
with drive $\drive(\ctrl) = \bm{\kappa}^{\top}\ctrl + \kappa_0$ an affine map
of the exogenous inputs only (injected power, current, mean density). For
$\lvert \drive \rvert < \drive_{\mathrm{f}} = 2 (b/3)^{3/2}$ the vector field
\eqref{eq:cusp} is bistable: stable equilibria
$\lat_{\mathrm{L}}^{*} < 0 < \lat_{\mathrm{H}}^{*}$ are separated by a saddle
$\lat_{\mathrm{s}}$. Sweeping $\drive$ through $+\drive_{\mathrm{f}}$
annihilates the L branch in a saddle-node bifurcation (forced L$\to$H
transition); sweeping back through $-\drive_{\mathrm{f}}$ annihilates the H
branch (H$\to$L back-transition). Hysteresis is therefore intrinsic rather
than imposed. We define the \emph{fold margins}
\begin{equation}
\mu_{\mathrm{LH}}(t) = \drive_{\mathrm{f}} - \drive(\ctrl(t)),
\qquad
\mu_{\mathrm{HL}}(t) = \drive(\ctrl(t)) + \drive_{\mathrm{f}},
\label{eq:margins}
\end{equation}
which measure the distance of the instantaneous drive from the two folds and
constitute the control-relevant output of the identified model.

\subsection{Observation model and regime probability}
The latent state is linked to a measured emission proxy $\emiss(t)$ (bright in
state L, quiescent in state H) through a learned map
$\hat{\emiss} = h_{\psi}\!\left(\barrier(\lat), \ctrl, \state_{\mathrm{edge}},
n_{\mathrm{edge}}\right)$. Crucially, $\emiss$ never enters the latent
dynamics \eqref{eq:cusp}; it only supervises the observation head, which
breaks the circularity of proxy-following formulations. The instantaneous
regime probability is the calibrated logistic
\begin{equation}
\pH(t) = \sigma\!\left(k\,\lat(t) / \sqrt{b}\right),
\label{eq:ph}
\end{equation}
and the mathematically sharp classification is basin membership:
the system is in regime H at time $t$ iff $\lat(t)$ lies in the basin of
attraction of $\lat^{*}_{\mathrm{H}}\big(\drive(t)\big)$.

\section{Identification}
\label{sec:ident}

\paragraph{Data.} Publicly released measurement campaigns from a toroidal
confinement experiment (MAST) \citep{mastopendata}; each discharge provides
profile measurements $\state(\rho_i, t_j)$ in the outer domain, exogenous
inputs $\ctrl(t)$, and the emission proxy $\emiss(t)$ at millisecond
resolution. The corpus contains 36 discharges: 26 with at least one L$\to$H
transition (several with back-transitions or dithering phases) and 10 that
remain in state L throughout, retained deliberately as negative examples.
Discharges are screened by strict per-channel quality gates (median level,
bounded jumps and span, minimum time coverage), which is essential: raw
channels are frequently corrupted, and admitting them destabilizes both the
rollout and the labels.

\paragraph{Weak segment labels.} Because $\emiss$ is bright in state L and
quiescent in state H but heavily contaminated by short bursts, labels are
computed from its \emph{lower envelope} (rolling 15th percentile, \SI{15}{ms}
window). Within a gate that excludes the formation and termination phases of
each discharge, the envelope distribution is split by a variance-maximizing
threshold; a candidate H segment is accepted only if
(i) it persists for at least \SI{20}{ms},
(ii) it is preceded by an L reference inside the gate, and
(iii) the envelope steps down sharply across its entry
(\SI{25}{ms} contrast windows), which rejects slow drifts.
Detected transition times agree with independent session-log annotations to
within $\sim$\SI{20}{ms} on the discharges where such annotations exist.

\paragraph{Objective.} With $\Theta = (\theta, \psi, \bm{\kappa}, b, \tau,
k, \chi\text{-parameters})$, the per-discharge loss is
\begin{equation}
\mathcal{L} =
\underbrace{\sum_{i,j} W_{ij}\, \phi_{\delta}\!\big(\hat\state(\rho_i, t_j) - \state_{ij}\big)}_{\text{profile misfit}}
+ \lambda_s \sum \phi_{\delta_s}(s_\theta)
+ \lambda_1 \overline{\lat^2} + \lambda_2 \overline{\dot\lat^2}
+ \lambda_r\, \mathrm{BCE}\big(\pH, \ell\big)
+ \lambda_y\, \overline{\big(h_\psi - \tilde\emiss\big)^2},
\label{eq:loss}
\end{equation}
where $\phi_\delta$ is the pseudo-Huber function, $W$ inverse-coverage
weights that equalize sparsely and densely observed channels, $\ell$ the weak
segment labels (scored only on clean L/H samples), and $\tilde\emiss$ the
range-normalized emission proxy. We use $\lambda_r = 5\times 10^{-2}$ and
$\lambda_y = 1$; in particular the emission term supervises the
\emph{observation head}, not the latent directly. Gradients of
\eqref{eq:loss} are obtained by reverse-mode differentiation through a
$\theta$-method implicit--explicit integrator ($\theta = 0.7$, five substeps
per observation interval) whose implicit stage is a tridiagonal (Thomas)
solve of the diffusion operator; the latent and source terms are explicit.
Model selection uses a held-out validation set of discharges whose loss
includes the regime and emission terms, so checkpoints are selected for
regime fidelity as well as profile fit.

\section{Results}
\label{sec:results}

\begin{figure}[t]
  \centering
  \includegraphics[width=0.72\linewidth]{figures/fig_data}
  \caption{Example discharge. Top: emission proxy $\emiss(t)$; the shaded band
  is the weakly labeled H segment and the dotted line the labeled transition
  time. Middle: exogenous inputs in fixed physical scales. Bottom: weak
  segment labels (L / transition / H; ``--'' marks the excluded formation and
  termination phases).}
  \label{fig:data}
\end{figure}

\paragraph{Identified normal form.} Trained on all 36 discharges (20\%
held out for model selection), the identified latent parameters are
$b = 1.00$, $\tau = \SI{23}{ms}$, and fold amplitude
$\drive_{\mathrm{f}} = 0.386$. In physical units the drive map gives a
forced-transition threshold ($\drive = +\drive_{\mathrm{f}}$) at
$\approx \SI{2.7}{MW}$ of injected power at the corpus-typical current and
density --- an interpretable, testable prediction of the identified model,
not an imposed constant. The relaxation time $\tau$ matches the observed
\SIrange{20}{40}{ms} rise of the transition, and the profile misfit in the
supervised outer annulus is \SI{73}{eV} (29\% relative).

\begin{figure}[t]
  \centering
  \includegraphics[width=0.6\linewidth]{figures/fig_bifurcation}
  \caption{Identified cusp equilibrium manifold (solid: stable branches;
  dashed: saddle) with the inferred trajectory
  $\big(\drive(\ctrl(t)), \lat(t)\big)$ of one discharge overlaid, colored by
  time. The state tracks the L branch, is carried past the fold
  $\drive_{\mathrm{f}}$ by the rising drive, and jumps to the H branch ---
  a saddle-node transition, not a fitted threshold.}
  \label{fig:bifurcation}
\end{figure}

\paragraph{Quantitative regime assessment.} Scoring the calibrated
probability $\pH(t)$ against the weak labels on clean L/H samples:
over the 26 discharges containing a labeled H segment, the median per-shot
AUC is $0.95$ (mean $0.85$ pooled over all 36), median accuracy $0.87$, and
the median error between the model's first sustained $\pH = 0.5$ upcrossing
and the labeled transition time is \SI{28}{ms} (mean \SI{45}{ms}).
All 10 non-transitioning discharges --- including purely resistively heated
ones --- are classified L throughout ($\pH$ never crosses $0.5$), which is
the falsifiability test that proxy-following formulations fail by
construction: here the latent is driven only by the actuators, yet it
declines to invent transitions where none occur.

\begin{figure}[t]
  \centering
  \includegraphics[width=0.66\linewidth]{figures/fig_classification}
  \caption{Regime probability $\pH(t)$ (blue) for four discharges; shaded
  bands are the weakly labeled H segments. The probability rises through the
  labeled transitions and remains near its branch value between them.}
  \label{fig:classification}
\end{figure}

\begin{figure}[t]
  \centering
  \includegraphics[width=0.68\linewidth]{figures/fig_margins}
  \caption{Fold margins along the discharge of Fig.~\ref{fig:bifurcation}.
  $\mu_{\mathrm{LH}} = \drive_{\mathrm{f}} - \drive(t)$ (blue) hits zero at
  the observed transition --- the L branch is annihilated --- while
  $\mu_{\mathrm{HL}} = \drive(t) + \drive_{\mathrm{f}}$ (orange) stays
  positive until the drive is removed, after which the H branch is lost and
  the state returns to L.}
  \label{fig:margins}
\end{figure}

\section{Discussion}
\label{sec:discussion}

The fold margins of Fig.~\ref{fig:margins} are the control-relevant output
of the identification: $\mu_{\mathrm{LH}}(t)$ is the instantaneous distance
(in drive units, hence convertible to watts) from a forced L$\to$H
transition, and $\mu_{\mathrm{HL}}(t)$ the distance from losing the H state.
A regulator that holds $\mu_{\mathrm{HL}}$ positive with a prescribed margin
maintains the high-confinement regime with a quantified robustness budget;
conversely, scheduling the drive through $+\drive_{\mathrm{f}}$ at a chosen
time forces the transition deliberately. Because both quantities are
closed-form functions of the identified parameters and the measured inputs,
they are computable in real time.

Limitations. First, the segment labels are weak: they derive from the same
emission proxy the observation head predicts (although they never enter the
latent dynamics), transition windows are short, and dithering phases are
labeled crudely; the classification ceiling is set by label quality.
Second, the latent-feature normalization of the legacy designs and the
per-shot standardization of the source-network inputs are non-causal
(whole-shot statistics); the cusp drive itself uses fixed physical scales
and is causal, but a fully real-time implementation must also make the
auxiliary inputs causal. Third, a single scalar mode cannot represent
oscillatory transition precursors (dithering limit cycles); an extension to
a two-dimensional normal form with a Hopf sector is natural future work.
Finally, the identified fold threshold is a property of this corpus
(one machine, one current/density neighborhood); transfer across operating
regimes requires letting $\drive(\ctrl)$ be nonlinear in the inputs, for
which the affine map here is the first term.

\bibliographystyle{plainnat}
\bibliography{refs}

\end{document}
